\documentclass[aspectratio=169]{beamer}

% Tema UFS --- Universidade Federal de Sergipe
\usetheme{UFS}

% Pacotes base
\usepackage[utf8]{inputenc}
\usepackage[portuguese]{babel}
\usepackage{amsmath,amssymb}
\usepackage{graphicx}
\usepackage{booktabs}
\usepackage{xcolor}

% TikZ e bibliotecas
\usepackage{tikz}
\usetikzlibrary{shapes.geometric, shapes.multipart, arrows.meta, positioning,
  matrix, fit, backgrounds, decorations.pathreplacing, patterns, calc}

% Listagens --- paleta VS Code Light+
\usepackage{listings}
\definecolor{bgcolor}{HTML}{FFFFFF}
\definecolor{linecolor}{HTML}{DDDDDD}
\definecolor{numcolor}{HTML}{999999}
\definecolor{kwcolor}{HTML}{0000FF}
\definecolor{strcolor}{HTML}{A31515}
\definecolor{cmtcolor}{HTML}{008000}

\lstdefinestyle{ufsStyle}{
  backgroundcolor=\color{bgcolor},
  basicstyle=\ttfamily\footnotesize\color{black},
  keywordstyle=\color{kwcolor}\bfseries,
  stringstyle=\color{strcolor},
  commentstyle=\color{cmtcolor}\itshape,
  numberstyle=\tiny\color{numcolor},
  numbers=left, numbersep=12pt,
  xleftmargin=20pt, framexleftmargin=20pt,
  breaklines=true, tabsize=4,
  keepspaces=true, showspaces=false,
  showstringspaces=false, showtabs=false,
  frame=single, rulecolor=\color{linecolor},
  framesep=4pt, captionpos=b, upquote=true,
  columns=fullflexible,
}
\lstset{style=ufsStyle, language=C}

\lstdefinestyle{tinyC}{
  style=ufsStyle, language=C, numbers=none,
  basicstyle=\ttfamily\scriptsize\color{black},
}
\lstdefinestyle{codC}{
  style=tinyC, basicstyle=\ttfamily\fontsize{7.5}{8.6}\selectfont\color{black},
}
\lstdefinestyle{shellStyle}{
  style=ufsStyle, language=bash, numbers=none,
  basicstyle=\ttfamily\scriptsize\color{black},
}

% --- Estilos TikZ reutilizados ---
\tikzset{
  cel/.style={rectangle, draw=UFSBlue, fill=UFSBlue!8, minimum width=0.55cm,
    minimum height=0.42cm, inner sep=0pt, font=\scriptsize\ttfamily},
  celvazia/.style={cel, fill=white, draw=UFSGray!60, text=UFSGray},
  celr/.style={cel, draw=UFSRed, fill=UFSRed!10},
  var/.style={rectangle, draw=UFSGray, fill=black!5, minimum width=0.7cm,
    minimum height=0.5cm, font=\scriptsize\ttfamily},
  varr/.style={var, draw=UFSRed, fill=UFSRed!8},
  leg/.style={font=\scriptsize, text=UFSGray},
  seta/.style={-{Stealth[length=1.8mm]}, thick, UFSBlue},
  setar/.style={-{Stealth[length=1.8mm]}, thick, UFSRed},
  setag/.style={-{Stealth[length=1.8mm]}, thick, UFSGray},
}

% Pilha desenhada de baixo para cima: \pilha{x}{y}{{a},{b},{c}} (c no topo)
\newcommand{\pilha}[3]{%
  \foreach \e [count=\k from 0] in {#3} {
    \node[cel] at (#1, #2+0.42*\k) {\e};
  }
  \draw[UFSGray, thick] (#1-0.36, #2+1.5) -- (#1-0.36, #2-0.24)
    -- (#1+0.36, #2-0.24) -- (#1+0.36, #2+1.5);
}
\newcommand{\pilhavazia}[2]{%
  \draw[UFSGray, thick] (#1-0.36, #2+1.5) -- (#1-0.36, #2-0.24)
    -- (#1+0.36, #2-0.24) -- (#1+0.36, #2+1.5);
}

% Referencia da fonte no rodape do frame
\newcommand{\fonte}[1]{\vfill\hfill{\scriptsize\textcolor{UFSGray}{Fonte: #1}}}

% --- Metadados ---
\title{Pilhas: conceito e implementação}
\subtitle{Só se põe e se tira por uma das pontas}
\author{Prof.\ André Yoshiaki Kashiwabara}
\institute{DCOMP -- Universidade Federal de Sergipe\\
\texttt{andreyoshiaki@dcomp.ufs.br}}
\date{}

\begin{document}

\begin{frame}[plain]
  \titlepage
\end{frame}

\begin{frame}{O que você vai aprender hoje}
  \begin{center}
  \begin{tikzpicture}[
    item/.style={rectangle, draw=UFSBlue, fill=UFSBlue!10, rounded corners,
      minimum width=10.4cm, minimum height=0.5cm, font=\small}
  ]
    \node[item] (t1) at (0, 0)    {Pilha: o último a entrar é o primeiro a sair};
    \node[item] (t2) at (0,-0.64) {Pilha num vetor: o topo é \texttt{v[n-1]}};
    \node[item] (t3) at (0,-1.28) {Pilha vazia, pilha cheia: os dois erros que não avisam};
    \node[item] (t4) at (0,-1.92) {Quando o vetor enche: dobrar, não somar um};
    \node[item] (t5) at (0,-2.56) {Pilha numa lista: o topo é o primeiro nó};
    \node[item] (t6) at (0,-3.20) {Quanto custa: vetor $\times$ lista};
    \node[item] (t7) at (0,-3.84) {Exercícios};
    \draw[seta] (t1) -- (t2);
    \draw[seta] (t2) -- (t3);
    \draw[seta] (t3) -- (t4);
    \draw[seta] (t4) -- (t5);
    \draw[seta] (t5) -- (t6);
    \draw[seta] (t6) -- (t7);
  \end{tikzpicture}
  \end{center}
\end{frame}

% ==================================================================
\section{A pergunta de hoje}
% ==================================================================

\begin{frame}{O problema: quem fecha quem?}
  \begin{columns}[T]
    \begin{column}{0.5\textwidth}
      \textbf{Até aqui}\\[0.3em]
      \small
      Com \texttt{ini} e \texttt{fim}, inserir e remover nas pontas da lista
      ficou barato. E na aula de recursão, a pilha de execução guardava cada
      chamada até ela terminar --- a última a começar era a primeira a acabar.
    \end{column}
    \begin{column}{0.5\textwidth}
      \textbf{O que vamos descobrir}\\[0.3em]
      \small
      Muitos problemas só precisam de \textbf{uma} ponta: guardar coisas e
      depois pegá-las na ordem inversa. Uma estrutura que só faz isso é mais
      simples --- e mais rápida --- do que uma lista.
    \end{column}
  \end{columns}
  \begin{block}{A pergunta de hoje}
    O compilador recusa \texttt{v[f(x)]} se faltar um \texttt{]}. Como um
    programa decide se \texttt{( [ \{ \} ] )} fecham na ordem certa --- e o
    que ele precisa guardar enquanto lê?
  \end{block}
\end{frame}

\begin{frame}[fragile]{O programa de hoje: \texttt{balanco.c} (1/2)}
  \begin{columns}[T]
    \begin{column}{0.46\textwidth}
\begin{lstlisting}[style=codC]
/* balanco.c --- ( [ { fecham
   na ordem certa? */
#include <stdio.h>

#define MAX 100

typedef struct {
    char v[MAX]; /* topo: v[n-1] */
    int n;       /* quantos ha */
} Pilha;

void empilha(Pilha *p, char c)
{
    p->v[p->n] = c;
    p->n++;
}
\end{lstlisting}
    \end{column}
    \begin{column}{0.5\textwidth}
\begin{lstlisting}[style=codC]
char desempilha(Pilha *p)
{
    p->n--;
    return p->v[p->n];
}

int vazia(const Pilha *p)
{
    return p->n == 0;
}

int combina(char abre, char fecha)
{
    return (abre == '(' && fecha == ')') ||
           (abre == '[' && fecha == ']') ||
           (abre == '{' && fecha == '}');
}
\end{lstlisting}
    \end{column}
  \end{columns}
\end{frame}

\begin{frame}[fragile]{\texttt{balanco.c} (2/2): a verificação e o \texttt{main}}
  \vspace{-0.4em}
  \begin{columns}[T]
    \begin{column}{0.54\textwidth}
\begin{lstlisting}[style=codC]
int balanceada(const char *s)
{
    Pilha p = {.n = 0};
    for (int i = 0; s[i] != '\0'; i++) {
        char c = s[i];
        if (c == '(' || c == '[' || c == '{')
            empilha(&p, c);
        else if (c == ')' || c == ']' || c == '}') {
            if (vazia(&p))
                return 0;  /* fechou sem abrir */
            if (!combina(desempilha(&p), c))
                return 0;  /* fechou o par errado */
        }
    }
    return vazia(&p);      /* sobrou aberto? */
}
\end{lstlisting}
    \end{column}
    \begin{column}{0.44\textwidth}
\begin{lstlisting}[style=codC]
int main(void)
{
    const char *testes[] = {
        "v[f(x)] = {1, 2}",
        "([)]",
        "(()",
        "())",
        "",
    };
    for (int i = 0; i < 5; i++)
        printf("%-18s %s\n", testes[i],
               balanceada(testes[i])
                   ? "sim" : "nao");
    return 0;
}
\end{lstlisting}
    \end{column}
  \end{columns}
\end{frame}

\begin{frame}{O que este programa imprime?}
  \begin{block}{Antes de compilar, escreva sua resposta}
    \begin{enumerate}
      \item Para cada uma das cinco cadeias, \texttt{sim} ou \texttt{nao}?
      \item Em \texttt{"([)]"}, o que está no topo da pilha quando chega o
        \texttt{)}?
      \item Em \texttt{"(()"} e em \texttt{"())"}, qual dos três
        \texttt{return} encerra a função?
      \item E a cadeia vazia, que não tem nada para fechar?
    \end{enumerate}
  \end{block}
\end{frame}

\begin{frame}[fragile]{Compile, execute e compare}
\begin{lstlisting}[style=shellStyle]
$ gcc -Wall -Wextra balanco.c -o balanco
$ ./balanco
v[f(x)] = {1, 2}   sim
([)]               nao
(()                nao
())                nao
                   sim
\end{lstlisting}
  \begin{columns}[T]
    \begin{column}{0.48\textwidth}
      \begin{block}{O que o programa guarda}
        \small Só os \textbf{abertos que ainda não fecharam} --- e o único
        que interessa a cada fechamento é o mais recente.
      \end{block}
    \end{column}
    \begin{column}{0.48\textwidth}
      \small Três maneiras de errar, três \texttt{return}: fechar o par
      errado (\texttt{([)]}), sobrar aberto (\texttt{(()}) e fechar sem ter
      aberto (\texttt{())}). A cadeia vazia é balanceada.
    \end{column}
  \end{columns}
\end{frame}

% ==================================================================
\section{O que é uma pilha}
% ==================================================================

\begin{frame}{Definição: pilha}
  \begin{columns}[c]
    \begin{column}{0.56\textwidth}
      \begin{block}{Pilha}
        \small Uma sequência em que inserções e remoções acontecem
        \textbf{só numa das pontas}, o \emph{topo}. O último a entrar é o
        primeiro a sair: \emph{LIFO}, \emph{last in, first out}.
      \end{block}
      \small
      \begin{tabular}{@{}ll@{}}
        \texttt{empilha(p, x)} & põe \texttt{x} no topo\\
        \texttt{desempilha(p)} & tira e devolve o topo\\
        \texttt{topo(p)}       & devolve o topo sem tirar\\
        \texttt{vazia(p)}      & a pilha não tem nada?
      \end{tabular}
    \end{column}
    \begin{column}{0.4\textwidth}
      \begin{center}
      \begin{tikzpicture}
        \pilha{0}{0}{{(},{[},{\{}}
        \node[leg, anchor=west] at (0.5,0.84) {$\leftarrow$ topo};
        \node[celr] (novo) at (-1.4,1.9) {\}};
        \draw[setar] (novo.east) to[out=0, in=90] (0,1.25);
        \node[leg] at (-1.4,1.5) {empilha};
        \draw[seta] (0.15,1.25) to[out=90, in=180] (1.2,1.9);
        \node[leg] at (1.5,1.5) {desempilha};
        \node[leg] at (0,-0.6) {fundo};
      \end{tikzpicture}
      \end{center}
    \end{column}
  \end{columns}
  \fonte{Cormen et al.\ (2009), §10.1~\cite{cormen2009introduction};
    Celes et al.\ (2004), cap.~11~\cite{celes2004introducao}}
\end{frame}

\begin{frame}{Seguindo a pilha caractere a caractere}
  \begin{center}
  \begin{tikzpicture}
    \node[leg, anchor=east] at (-0.8,0.5) {\texttt{([)]}};
    \pilha{0}{0}{{(}}         \node[leg] at (0,-0.5) {\texttt{(}};
    \pilha{1.2}{0}{{(},{[}}     \node[leg] at (1.2,-0.5) {\texttt{[}};
    \node[celr] at (2.4,0.42) {[};
    \pilha{2.4}{0}{{(}}
    \node[leg] at (2.4,-0.5) {\texttt{)}};
    \node[leg, text=UFSRed, anchor=west] at (2.9,0.42)
      {sai \texttt{[}, chegou \texttt{)}: não combina};

    \node[leg, anchor=east] at (-0.8,-1.7) {\texttt{v[f(x)]}};
    \pilha{0}{-2.2}{{[}}      \node[leg] at (0,-2.7) {\texttt{[}};
    \pilha{1.2}{-2.2}{{[},{(}}  \node[leg] at (1.2,-2.7) {\texttt{(}};
    \pilha{2.4}{-2.2}{{[}}    \node[leg] at (2.4,-2.7) {\texttt{)}};
    \pilhavazia{3.6}{-2.2}  \node[leg] at (3.6,-2.7) {\texttt{]}};
    \node[leg, anchor=west] at (4.1,-1.8) {letras e \texttt{=} passam direto;};
    \node[leg, anchor=west] at (4.1,-2.15) {no fim, pilha vazia: \texttt{sim}};
  \end{tikzpicture}
  \end{center}
  \small Cada fechamento só confere o \textbf{topo}: o aberto mais recente
  ainda sem par. Os de baixo esperam a vez --- nenhum precisa ser procurado.
  \fonte{Tenenbaum et al.\ (1995), cap.~2~\cite{tenenbaum1995estruturas};
    Feofiloff (2009), cap.~6~\cite{feofiloff2009algoritmos}}
\end{frame}

\begin{frame}{Pilhas que você já usou}
  \begin{columns}[T]
    \begin{column}{0.32\textwidth}
      \begin{block}{Pilha de execução}
        \small Cada chamada empilha um registro de ativação; cada
        \texttt{return} desempilha. Foi assim que \texttt{fatorial(4)}
        voltou na ordem inversa.
      \end{block}
    \end{column}
    \begin{column}{0.32\textwidth}
      \begin{block}{Desfazer}
        \small Cada edição vai para uma pilha; \textsf{Ctrl+Z} tira a mais
        recente. Desfazer duas vezes desfaz as duas últimas, na ordem
        inversa.
      \end{block}
    \end{column}
    \begin{column}{0.32\textwidth}
      \begin{block}{Voltar no navegador}
        \small Cada página visitada é empilhada; ``voltar'' desempilha.
        Abrir um link novo depois de voltar descarta o que estava acima.
      \end{block}
    \end{column}
  \end{columns}
  \medskip
  \small Em todos, a pergunta é a mesma: \emph{qual foi a última coisa que
  ainda não terminou?}
  \fonte{Knuth (1997), §2.2.1~\cite{knuth1997art};
    Tenenbaum et al.\ (1995), cap.~2~\cite{tenenbaum1995estruturas}}
\end{frame}

\begin{frame}{Recapitulando: o que é uma pilha}
  \begin{block}{O que aprendemos}
    \begin{itemize}
      \item pilha: inserir e remover só no topo --- o último a entrar é o
        primeiro a sair;
      \item quatro operações: \texttt{empilha}, \texttt{desempilha},
        \texttt{topo} e \texttt{vazia};
      \item para balancear, basta guardar os abertos ainda sem par: cada
        fechamento confere só o topo;
      \item a pilha de execução da recursão é uma pilha no mesmo sentido.
    \end{itemize}
  \end{block}
\end{frame}

% ==================================================================
\section{Pilha num vetor}
% ==================================================================

\begin{frame}{O topo é \texttt{v[n-1]}}
  \begin{center}
  \begin{tikzpicture}
    \foreach \c [count=\k from 0] in {{(},{[},{\{}} {
      \node[cel, minimum width=0.7cm] (v\k) at (0.7*\k,0) {\c};
      \node[leg] at (0.7*\k,-0.4) {\k};
    }
    \foreach \k in {3,...,7} {
      \node[celvazia, minimum width=0.7cm] (v\k) at (0.7*\k,0) {?};
      \node[leg] at (0.7*\k,-0.4) {\k};
    }
    \node[var] (n) at (-1.6,0) {n = 3};
    \draw[setar] (v2.north) ++(0,0.5) node[above left=0pt and -6pt, leg, text=UFSRed] {topo: \texttt{v[n-1]}} -- (v2.north);
    \draw[seta] (v3.north) ++(0,0.5) node[above right=0pt and -6pt, leg, text=UFSBlue] {próximo: \texttt{v[n]}} -- (v3.north);
    \node[leg, anchor=west] at (5.3,0) {\ldots\ \texttt{MAX-1}};
  \end{tikzpicture}
  \end{center}
  \begin{columns}[T]
    \begin{column}{0.48\textwidth}
      \small
      \texttt{empilha}: escreve em \texttt{v[n]}, depois \texttt{n++}.\\
      \texttt{desempilha}: \texttt{n-{}-}, depois lê \texttt{v[n]}.\\
      \texttt{vazia}: \texttt{n == 0}.
    \end{column}
    \begin{column}{0.48\textwidth}
      \small Nenhuma operação tem laço: cada uma custa um número fixo de
      passos. O fundo nunca se move --- é o topo que anda. Desempilhar não
      apaga nada: o valor continua lá, fora da pilha.
    \end{column}
  \end{columns}
  \fonte{Cormen et al.\ (2009), §10.1, \textsc{Push} e \textsc{Pop}~\cite{cormen2009introduction};
    Celes et al.\ (2004), cap.~11~\cite{celes2004introducao}}
\end{frame}

\begin{frame}[fragile]{E se não conferir a pilha vazia?}
  \begin{columns}[c]
    \begin{column}{0.54\textwidth}
\begin{lstlisting}[style=tinyC]
/* balanco_errada.c: o laco de balanceada */
    for (int i = 0; s[i] != '\0'; i++) {
        char c = s[i];
        if (c == '(' || c == '[' || c == '{')
            empilha(&p, c);
        else if (c == ')' || c == ']' || c == '}') {
            /* sumiu: if (vazia(&p)) return 0; */
            if (!combina(desempilha(&p), c))
                return 0;
        }
    }
\end{lstlisting}
    \end{column}
    \begin{column}{0.42\textwidth}
      \begin{block}{Pense antes de passar o slide}
        \small O \texttt{main} é o mesmo.
        \begin{enumerate}
          \item Qual das cinco cadeias chega a desempilhar a pilha vazia?
          \item Quanto vale \texttt{p->n} depois, e que posição de
            \texttt{v} é lida?
          \item A resposta impressa muda?
        \end{enumerate}
      \end{block}
    \end{column}
  \end{columns}
\end{frame}

\begin{frame}[fragile]{Execute e compare}
\begin{lstlisting}[style=shellStyle]
$ ./balanco_errada
v[f(x)] = {1, 2}   sim
([)]               nao
(()                nao
())                nao
                   sim
$ gcc -g -fsanitize=address balanco_errada.c -o balanco_errada-asan
$ ./balanco_errada-asan
==95953==ERROR: AddressSanitizer: stack-buffer-underflow on address 0x00016f39245f
READ of size 1 at 0x00016f39245f thread T0
    #0 0x000100a6cab4 in desempilha balanco_errada.c:20
    #1 0x000100a6cea0 in balanceada balanco_errada.c:44
  This frame has 1 object(s):
    [32, 136) 'p' (line 38) <== Memory access at offset 31 underflows this variable
\end{lstlisting}
  \begin{center}
    \scriptsize\textcolor{UFSGray}{A linha 20 é o \texttt{return p->v[p->n];}
    de \texttt{desempilha}. Saída do relatório resumida.}
  \end{center}
\end{frame}

\begin{frame}{A resposta: certo por sorte é errado}
  \begin{columns}[c]
    \begin{column}{0.48\textwidth}
      \begin{center}
      \begin{tikzpicture}
        \node[celr, minimum width=0.7cm] (m1) at (0,0) {??};
        \node[leg] at (0,-0.4) {-1};
        \foreach \k in {0,...,3} {
          \node[celvazia, minimum width=0.7cm] at (0.7+0.7*\k,0) {};
          \node[leg] at (0.7+0.7*\k,-0.4) {\k};
        }
        \draw[UFSBlue, thick, decorate, decoration={brace, amplitude=4pt}]
          (0.35,0.3) -- (3.15,0.3) node[midway, above=3pt, leg, text=UFSBlue] {\texttt{p.v}};
        \node[var] at (1.75,-1.1) {n: 0 $\to$ -1};
        \draw[setar] (0,1.0) node[above, leg, text=UFSRed] {\texttt{v[-1]}: fora da pilha} -- (m1.north);
      \end{tikzpicture}
      \end{center}
    \end{column}
    \begin{column}{0.48\textwidth}
      \small
      Em \texttt{"())"}, o segundo \texttt{)} encontra a pilha vazia:
      \texttt{n} vira $-1$ e \texttt{desempilha} lê \texttt{v[-1]}, um byte
      \textbf{antes} do vetor.

      \medskip
      O byte lido era lixo, \texttt{combina} devolveu 0 e a resposta saiu
      \texttt{nao} --- a certa, por acaso. Ler fora do vetor é comportamento
      indefinido: outro compilador, outra resposta.

      \medskip
      Desempilhar a pilha vazia tem nome: \emph{underflow}. Quem chama
      precisa perguntar \texttt{vazia} antes.
    \end{column}
  \end{columns}
  \fonte{Cormen et al.\ (2009), §10.1, \emph{underflow}~\cite{cormen2009introduction};
    ISO/IEC 9899:2018, §6.5.6~\cite{iso2018c17}; Serebryany et al.\ (2012)~\cite{serebryany2012asan}}
\end{frame}

\begin{frame}[fragile]{E se a pilha encher?}
  \begin{columns}[c]
    \begin{column}{0.54\textwidth}
\begin{lstlisting}[style=tinyC]
/* transborda.c: k vezes '(', depois ')' */
int niveis[] = {99, 100, 101, 150};
for (int i = 0; i < 4; i++) {
    Pilha p = {.n = 0};
    int k = niveis[i], fechou = 0;
    for (int j = 0; j < k; j++)
        empilha(&p, '(');
    printf("%3d abertos: n = %3d", k, p.n);
    while (!vazia(&p) && fechou < k) {
        desempilha(&p);
        fechou++;
    }
    int ok = fechou == k && vazia(&p);
    printf(", fechou %3d -> %s\n", fechou,
           ok ? "sim" : "nao");
}
\end{lstlisting}
    \end{column}
    \begin{column}{0.42\textwidth}
      \begin{block}{Pense antes de passar o slide}
        \small \texttt{MAX} continua 100.
        \begin{enumerate}
          \item Para \texttt{k = 101}, onde o último \texttt{(} é escrito?
          \item Que valor de \texttt{n} você espera ver impresso?
          \item O \texttt{-fsanitize=address} vai reclamar?
        \end{enumerate}
      \end{block}
    \end{column}
  \end{columns}
\end{frame}

\begin{frame}[fragile]{Execute e compare}
\begin{lstlisting}[style=shellStyle]
$ ./transborda
 99 abertos: n =  99, fechou  99 -> sim
100 abertos: n = 100, fechou 100 -> sim
101 abertos: n =  41, fechou  41 -> nao
150 abertos: n =  90, fechou  90 -> nao
$ gcc -g -fsanitize=address transborda.c -o transborda-asan && ./transborda-asan
 99 abertos: n =  99, fechou  99 -> sim
100 abertos: n = 100, fechou 100 -> sim
101 abertos: n =  41, fechou  41 -> nao
150 abertos: n =  90, fechou  90 -> nao
\end{lstlisting}
  \begin{columns}[T]
    \begin{column}{0.48\textwidth}
      \small Empilhou 101, mas \texttt{n} diz 41. E nem o
      \emph{sanitizer} viu nada.
    \end{column}
    \begin{column}{0.48\textwidth}
      \small Para 150, \texttt{n} termina em 90: os 49 \texttt{(} depois do
      estrago foram escritos \emph{dentro} do vetor, por cima dos antigos.
    \end{column}
  \end{columns}
\end{frame}

\begin{frame}{A resposta: \texttt{v[100]} é o primeiro byte de \texttt{n}}
  \begin{center}
  \begin{tikzpicture}
    \node[cel, minimum width=0.7cm] (a) at (0,0) {(};
    \node[leg] at (0,-0.4) {0};
    \node[leg] at (0.75,0) {\ldots};
    \node[cel, minimum width=0.7cm] (b) at (1.5,0) {(};
    \node[leg] at (1.5,-0.4) {99};
    \node[celr, minimum width=0.7cm] (c) at (2.2,0) {40};
    \foreach \k in {1,2,3} \node[cel, fill=black!5, draw=UFSGray, minimum width=0.7cm] at (2.2+0.7*\k,0) {0};
    \draw[UFSBlue, thick, decorate, decoration={brace, amplitude=4pt}]
      (-0.35,0.3) -- (1.85,0.3) node[midway, above=3pt, leg, text=UFSBlue] {\texttt{v}: 100 bytes};
    \draw[UFSGray, thick, decorate, decoration={brace, amplitude=4pt}]
      (1.85,0.3) -- (4.65,0.3) node[midway, above=3pt, leg] {\texttt{n}: 4 bytes};
    \draw[setar] (2.2,-1.0) node[below, leg, text=UFSRed] {\texttt{v[100] = '('} grava 40} -- (c.south);
  \end{tikzpicture}
  \end{center}
  \small
  No registro, \texttt{n} vem logo depois dos 100 bytes de \texttt{v}. O 101º
  \texttt{empilha} escreve \texttt{'('} --- código 40 --- no byte menos
  significativo de \texttt{n}, que valia 100 (x86-64 e ARM guardam esse byte
  primeiro); o \texttt{n++} seguinte faz 41. A escrita caiu dentro do objeto
  \texttt{p}, e o \emph{sanitizer} não a acusa. Encher a pilha é o \emph{overflow}: quem empilha precisa
  saber se ainda cabe.
  \fonte{ISO/IEC 9899:2018, §6.7.2.1 e §6.5.6~\cite{iso2018c17};
    Cormen et al.\ (2009), §10.1~\cite{cormen2009introduction}; medido com \texttt{transborda.c}}
\end{frame}

\begin{frame}{Recapitulando: pilha num vetor}
  \begin{block}{O que aprendemos}
    \begin{itemize}
      \item o fundo é \texttt{v[0]}, o topo é \texttt{v[n-1]}; empilhar e
        desempilhar só mexem em \texttt{n} e numa posição;
      \item desempilhar a vazia (\emph{underflow}) lê fora do vetor;
        empilhar na cheia (\emph{overflow}) escreve fora;
      \item os dois podem dar a resposta certa, uma errada ou nenhum aviso
        --- nem do ASan;
      \item conferir \texttt{vazia} antes de tirar e ``ainda cabe?'' antes de
        pôr não é opcional.
    \end{itemize}
  \end{block}
\end{frame}

% ==================================================================
\section{Quando o vetor enche}
% ==================================================================

\begin{frame}[fragile]{Um vetor que cresce: \texttt{pilha\_vetor.c}}
  \begin{columns}[T]
    \begin{column}{0.5\textwidth}
\begin{lstlisting}[style=codC]
struct pilha {
    int *v;    /* topo: v[n-1] */
    int n;     /* quantos ha */
    int cap;   /* quantos cabem */
};

void empilha(Pilha *p, int x)
{
    if (p->n == p->cap) { /* cheio */
        p->cap *= 2;
        p->v = aloca(p->v,
                 p->cap * sizeof *p->v);
    }
    p->v[p->n] = x;
    p->n++;
}
\end{lstlisting}
    \end{column}
    \begin{column}{0.46\textwidth}
      \small \texttt{aloca} é um \texttt{realloc} que encerra o programa se
      faltar memória; \texttt{pilha\_cria} começa com \texttt{cap = 4}.
      \begin{block}{Pense antes de passar o slide}
        \small
        \begin{enumerate}
          \item Por que não \texttt{p->cap += 1}? Gasta menos memória.
          \item Para empilhar $n$ valores a partir de \texttt{cap = 1},
            quantos elementos são copiados em cada estratégia, se cada
            \texttt{realloc} mudar o bloco de lugar?
        \end{enumerate}
      \end{block}
    \end{column}
  \end{columns}
  \fonte{ISO/IEC 9899:2018, §7.22.3.5~\cite{iso2018c17}}
\end{frame}

\begin{frame}{Execute e compare: somar um $\times$ dobrar}
  \begin{columns}[c]
    \begin{column}{0.56\textwidth}
      \begin{center}
      \small
      \begin{tabular}{@{}rrrrr@{}}
        \toprule
        & \multicolumn{2}{c}{\texttt{cap + 1}} & \multicolumn{2}{c}{\texttt{cap * 2}}\\
        \cmidrule(lr){2-3}\cmidrule(l){4-5}
        $n$ & cópias & tempo & cópias & tempo\\
        \midrule
        20\,000  & $2{,}0\cdot 10^{8}$ & 0,05\,s & 32\,767  & $<$0,001\,s\\
        40\,000  & $8{,}0\cdot 10^{8}$ & 0,34\,s & 65\,535  & $<$0,001\,s\\
        80\,000  & $3{,}2\cdot 10^{9}$ & 0,61\,s & 131\,071 & $<$0,001\,s\\
        160\,000 & $1{,}3\cdot 10^{10}$& 1,65\,s & 262\,143 & $<$0,001\,s\\
        \bottomrule
      \end{tabular}

      \medskip
      \scriptsize \texttt{crescimento.c}, \texttt{-O2}; cada crescimento
      aloca, copia e libera
      \end{center}
    \end{column}
    \begin{column}{0.4\textwidth}
      \small
      Somando um, dobrar $n$ \textbf{quadruplica} as cópias.

      \medskip
      Dobrando, as cópias ficam sempre abaixo de $2n$: dobrar $n$ dobra
      as cópias.

      \medskip
      \scriptsize O \texttt{realloc} de verdade às vezes aumenta o bloco no
      lugar e não copia nada --- mas a norma não promete isso.
    \end{column}
  \end{columns}
  \fonte{medido com \texttt{crescimento.c}; ISO/IEC 9899:2018, §7.22.3.5~\cite{iso2018c17}}
\end{frame}

\begin{frame}{A resposta: dobrar paga as cópias adiantado}
  \begin{columns}[c]
    \begin{column}{0.5\textwidth}
      \begin{center}
      \begin{tikzpicture}[x=0.32cm, y=0.32cm]
        \foreach \c/\y in {1/0, 2/-1.3, 4/-2.6, 8/-3.9, 16/-5.2} {
          \foreach \i in {1,...,\c} {
            \pgfmathtruncatemacro{\cheio}{\i <= \c/2 ? 1 : 0}
            \ifnum\cheio=1
              \draw[UFSRed, fill=UFSRed!25] (\i-1,\y) rectangle ++(1,1);
            \else
              \draw[UFSBlue, fill=UFSBlue!8] (\i-1,\y) rectangle ++(1,1);
            \fi
          }
          \node[leg, anchor=east] at (-0.3,\y+0.5) {\c};
        }
        \node[leg, anchor=west] at (0,-6.4) {\textcolor{UFSRed}{vermelho}: copiado ao crescer};
      \end{tikzpicture}
      \end{center}
    \end{column}
    \begin{column}{0.46\textwidth}
      \small
      \textbf{Somando um}: $1 + 2 + \cdots + (n-1) = n(n-1)/2$ cópias.\\[0.4em]
      \textbf{Dobrando}: $1 + 2 + 4 + \cdots + n/2 < n$ cópias.\\[0.6em]
      Cada crescimento copia metade do novo vetor, mas depois dele cabem
      outros tantos \texttt{empilha} sem cópia nenhuma. Dividido entre todos,
      o custo por \texttt{empilha} é constante: o \emph{custo amortizado}.

      \medskip
      O preço: até metade do vetor pode ficar vazia.
    \end{column}
  \end{columns}
  \fonte{Cormen et al.\ (2009), §17.4, \emph{dynamic tables}~\cite{cormen2009introduction}}
\end{frame}

\begin{frame}{Recapitulando: quando o vetor enche}
  \begin{block}{O que aprendemos}
    \begin{itemize}
      \item guardar \texttt{cap} ao lado de \texttt{n} e crescer com
        \texttt{realloc} quando \texttt{n == cap} acaba com o \emph{overflow};
      \item crescer de um em um custa $n^2/2$ cópias; dobrar, menos de $n$;
      \item dobrando, um \texttt{empilha} raro é caro, mas a média por
        operação é constante;
      \item a troca: até metade da memória reservada pode ficar sem uso.
    \end{itemize}
  \end{block}
\end{frame}

% ==================================================================
\section{Pilha numa lista}
% ==================================================================

\begin{frame}[fragile]{O topo é o primeiro nó: \texttt{pilha\_lista.c}}
  \begin{columns}[T]
    \begin{column}{0.5\textwidth}
\begin{lstlisting}[style=codC]
struct pilha {
    No *topo;   /* 1o no, ou NULL */
};

void empilha(Pilha *p, int x)
{   /* insere no inicio */
    No *novo = malloc(sizeof *novo);
    if (novo == NULL) exit(1);
    novo->valor = x;
    novo->prox = p->topo;
    p->topo = novo;
}
\end{lstlisting}
    \end{column}
    \begin{column}{0.46\textwidth}
\begin{lstlisting}[style=codC]
int desempilha(Pilha *p)
{   /* remove o primeiro */
    if (p->topo == NULL) {
        fprintf(stderr, "pilha vazia\n");
        exit(1);
    }
    No *velho = p->topo;
    int x = velho->valor;
    p->topo = velho->prox;
    free(velho);
    return x;
}
\end{lstlisting}
    \end{column}
  \end{columns}
  \begin{center}
  \begin{tikzpicture}
    \node[var] (t) at (0,0) {topo};
    \foreach \v/\k in {{\{}/1, {[}/2, {(}/3} {
      \node[cel, minimum width=0.6cm] (v\k) at (0.7+1.3*\k,0) {\v};
      \node[cel, fill=white, minimum width=0.3cm, right=0pt of v\k] (p\k) {};
    }
    \draw[seta] (t) -- (v1);
    \draw[seta] (p1.center) -- (v2); \draw[seta] (p2.center) -- (v3);
    \draw[UFSGray] (p3.south west) -- (p3.north east);
    \node[leg, anchor=west] at (5.55,0) {o \texttt{(} é o fundo: foi o primeiro a entrar};
  \end{tikzpicture}
  \end{center}
  \fonte{Celes et al.\ (2004), cap.~11~\cite{celes2004introducao};
    Feofiloff (2009), cap.~6~\cite{feofiloff2009algoritmos}}
\end{frame}

\begin{frame}{Por que a lista simples basta}
  \begin{columns}[T]
    \begin{column}{0.48\textwidth}
      \small
      Inserir e remover no \textbf{início} eram as duas operações baratas da
      lista simples: nenhuma precisa de \texttt{ant}, de \texttt{fim} ou de
      laço.

      \medskip
      Se o topo fosse o \emph{último} nó, empilhar seria barato com
      \texttt{fim}, mas desempilhar exigiria achar o penúltimo --- $n$
      passos, ou um \texttt{ant} em cada nó.
    \end{column}
    \begin{column}{0.48\textwidth}
      \begin{block}{Sem \emph{overflow}, sem \texttt{cap}}
        \small Cada \texttt{empilha} pede um nó ao \texttt{malloc}: a pilha
        só enche quando acaba a memória do processo. O \emph{underflow}
        continua possível --- por isso o teste em \texttt{desempilha}.
      \end{block}
    \end{column}
  \end{columns}
  \fonte{Celes et al.\ (2004), cap.~11~\cite{celes2004introducao};
    Cormen et al.\ (2009), §10.2, exercício 10.2-2~\cite{cormen2009introduction}}
\end{frame}

\begin{frame}[fragile]{Uma interface, duas implementações}
  \begin{columns}[T]
    \begin{column}{0.5\textwidth}
\begin{lstlisting}[style=codC]
/* pilha.h */
#ifndef PILHA_H
#define PILHA_H

typedef struct pilha Pilha; /* campos no .c */

Pilha *pilha_cria(void);
void pilha_libera(Pilha *p);

void empilha(Pilha *p, int x);
int desempilha(Pilha *p);
int topo(const Pilha *p);
int vazia(const Pilha *p);

#endif
\end{lstlisting}
    \end{column}
    \begin{column}{0.46\textwidth}
      \small \texttt{balanco\_tad.c} tem a mesma \texttt{balanceada}, agora
      com \texttt{Pilha *}, e testa também $10^5$ abertos seguidos de $10^5$
      fechados.
      \begin{block}{Pense antes de passar o slide}
        \small
        \begin{enumerate}
          \item O \texttt{main} pode escrever \texttt{p->n}? E
            \texttt{sizeof(Pilha)}?
          \item O que muda em \texttt{balanco\_tad.c} para trocar o vetor
            pela lista?
          \item As duas versões imprimem \texttt{sim} para $10^5$ níveis?
        \end{enumerate}
      \end{block}
    \end{column}
  \end{columns}
  \fonte{Kernighan e Ritchie (1988), §4.5~\cite{kernighan1988c};
    Celes et al.\ (2004), cap.~9~\cite{celes2004introducao}}
\end{frame}

\begin{frame}[fragile]{Execute e compare}
\begin{lstlisting}[style=shellStyle]
$ gcc -Wall -Wextra balanco_tad.c pilha_vetor.c -o balanco_vetor
$ gcc -Wall -Wextra balanco_tad.c pilha_lista.c -o balanco_lista
$ ./balanco_vetor                     $ ./balanco_lista
v[f(x)] = {1, 2}   sim                v[f(x)] = {1, 2}   sim
([)]               nao                ([)]               nao
(()                nao                (()                nao
())                nao                ())                nao
                   sim                                   sim
100000 niveis      sim                100000 niveis      sim
\end{lstlisting}
  \begin{columns}[T]
    \begin{column}{0.48\textwidth}
      \small \texttt{balanco\_tad.c} não mudou \textbf{nenhuma} linha: só
      o arquivo passado ao \texttt{gcc}. E $10^5$ níveis já não transbordam.
    \end{column}
    \begin{column}{0.48\textwidth}
      \small O \texttt{main} só vê um \texttt{struct pilha} incompleto: pode
      guardar um \texttt{Pilha *}, mas \texttt{p->n} e
      \texttt{sizeof(Pilha)} não compilam.
    \end{column}
  \end{columns}
\end{frame}

\begin{frame}{Recapitulando: pilha numa lista}
  \begin{block}{O que aprendemos}
    \begin{itemize}
      \item o topo é o primeiro nó: empilhar é inserir no início,
        desempilhar é remover o primeiro;
      \item uma lista simples basta --- só uma ponta é usada;
      \item a lista não enche, mas desempilhar a vazia continua sendo erro;
      \item com \texttt{pilha.h}, quem usa a pilha não sabe nem precisa
        saber se ela é vetor ou lista.
    \end{itemize}
  \end{block}
\end{frame}

% ==================================================================
\section{Quanto custa}
% ==================================================================

\begin{frame}{Empilhar e desempilhar $n$ valores: vetor $\times$ lista}
  \begin{columns}[c]
    \begin{column}{0.48\textwidth}
      \begin{center}
      \small
      \begin{tabular}{@{}rrr@{}}
        \toprule
        $n$ & vetor (dobra) & lista\\
        \midrule
        1\,000\,000 & 0,002\,s & 0,035\,s\\
        2\,000\,000 & 0,003\,s & 0,123\,s\\
        4\,000\,000 & 0,005\,s & 0,127\,s\\
        8\,000\,000 & 0,009\,s & 0,234\,s\\
        \bottomrule
      \end{tabular}

      \medskip
      \scriptsize \texttt{vetor\_vs\_lista.c}, \texttt{-O2}, valores
      \texttt{int}
      \end{center}
    \end{column}
    \begin{column}{0.48\textwidth}
      \small
      As duas crescem em linha reta: cada operação custa um número fixo de
      passos. Mas a lista é umas \textbf{25 vezes} mais lenta.

      \medskip
      Na lista, cada \texttt{empilha} é um \texttt{malloc} e cada
      \texttt{desempilha}, um \texttt{free}; os nós ficam espalhados pela
      memória. No vetor, os valores estão lado a lado e o
      \texttt{realloc} é raro.
    \end{column}
  \end{columns}
  \fonte{medido com \texttt{vetor\_vs\_lista.c};
    Bryant e O'Hallaron (2016), §6.2 e §9.9~\cite{bryant2016computer}}
\end{frame}

\begin{frame}{Vetor ou lista?}
  \begin{center}
  \small
  \begin{tabular}{@{}lccc@{}}
    \toprule
    & vetor fixo & vetor que dobra & lista encadeada\\
    \midrule
    \texttt{empilha}             & constante & constante (amortizado) & constante + \texttt{malloc}\\
    \texttt{desempilha}, \texttt{topo}, \texttt{vazia} & constante & constante & constante\\
    pode encher?                 & sim, em \texttt{MAX} & só sem memória & só sem memória\\
    memória por \texttt{int}     & \texttt{MAX} $\times$ 4\,B & 4 a 8\,B & 16\,B + o do \texttt{malloc}\\
    \bottomrule
  \end{tabular}
  \end{center}
  \small Para uma pilha, o vetor que dobra costuma ganhar. A lista compensa
  quando cada elemento é grande e já mora num nó --- ou quando um
  \texttt{empilha} nunca pode demorar mais que os outros.
  \fonte{Cormen et al.\ (2009), §10.1 e §17.4~\cite{cormen2009introduction};
    Celes et al.\ (2004), cap.~11~\cite{celes2004introducao}}
\end{frame}

\begin{frame}{Recapitulando: quanto custa}
  \begin{block}{O que aprendemos}
    \begin{itemize}
      \item nas três versões, toda operação de pilha tem custo fixo --- não
        depende de $n$;
      \item o vetor que dobra é o mais rápido na prática: memória contígua e
        poucas chamadas ao alocador;
      \item a lista gasta 16 bytes por \texttt{int} e um \texttt{malloc} por
        \texttt{empilha};
      \item o vetor fixo só serve quando há um limite conhecido --- e
        conferido.
    \end{itemize}
  \end{block}
\end{frame}

% ==================================================================
\section{Mãos à obra}
% ==================================================================

\begin{frame}{Altere o programa}
  \small
  \begin{block}{Tarefa 1}
    Faça \texttt{balanceada} devolver a \textbf{posição} do primeiro erro (ou
    $-1$ se estiver tudo certo) e imprima-a. Em \texttt{"(()"}, qual posição
    faz sentido informar?
  \end{block}
  \begin{block}{Tarefa 2}
    Guarde na pilha a posição de cada aberto (um \texttt{int}), em vez do
    caractere. O que \texttt{combina} passa a receber? Que mensagem fica
    possível em \texttt{"(()"}?
  \end{block}
  \begin{block}{Tarefa 3}
    Troque \texttt{MAX} por 4 e acrescente a \texttt{empilha} um teste que
    encerra com mensagem se a pilha estiver cheia. Rode
    \texttt{transborda.c} de novo e preveja cada linha.
  \end{block}
\end{frame}

\begin{frame}{Exercícios (1/3)}
  \footnotesize
  \begin{block}{Exercício 1 --- uma pilha que avisa}
    \texttt{int empilha(Pilha *p, char c)} devolve 0 se a pilha está cheia;
    \texttt{int desempilha(Pilha *p, char *c)} devolve 0 se está vazia e o
    topo por \texttt{*c}. Use-as em \texttt{balanceada}.
  \end{block}
  \begin{block}{Exercício 2 --- inverter uma cadeia}
    \texttt{void inverte(char *s)} com uma pilha de \texttt{char}; depois
    \texttt{inverte\_palavras}: \texttt{"ola mundo"} $\to$ \texttt{"alo odnum"}.
  \end{block}
  \begin{block}{Exercício 3 --- decimal para binário}
    Imprima os dígitos de \texttt{n} na base \texttt{b} ($2 \le b \le 16$):
    os restos de \texttt{n \% b} saem do menos significativo para o mais ---
    a pilha devolve na ordem certa.
  \end{block}
  \fonte{Celes et al.\ (2004), cap.~11~\cite{celes2004introducao};
    Tenenbaum et al.\ (1995), cap.~2~\cite{tenenbaum1995estruturas}}
\end{frame}

\begin{frame}{Exercícios (2/3)}
  \footnotesize
  \begin{block}{Exercício 4 --- calculadora pós-fixa}
    Avalie \texttt{"3 4 + 2 *"} $\to$ 14 (pós-fixa, \texttt{+ - * /}).
    Número: empilha. Operador: desempilha dois, empilha o resultado. Qual
    dos dois é o operando da esquerda?
  \end{block}
  \begin{block}{Exercício 5 --- pilha com mínimo}
    Acrescente \texttt{int minimo(const Pilha *p)}, que devolve o menor
    elemento em tempo constante. Dica: uma segunda pilha de mínimos.
  \end{block}
  \begin{block}{Exercício 6 --- duas pilhas num vetor}
    Num único \texttt{int v[MAX]}, uma pilha cresce de \texttt{v[0]} para a
    direita e outra de \texttt{v[MAX-1]} para a esquerda. Quando enchem?
  \end{block}
  \fonte{Feofiloff (2009), cap.~6~\cite{feofiloff2009algoritmos};
    Knuth (1997), §2.2.2~\cite{knuth1997art}}
\end{frame}

\begin{frame}{Exercícios (3/3)}
  \footnotesize
  \begin{block}{Exercício 7 --- desfazer e refazer}
    Um contador recebe comandos \texttt{+k}, \texttt{-k}, \texttt{u}
    (desfaz) e \texttt{r} (refaz). Use duas pilhas. O que acontece com a
    pilha de refazer quando chega um \texttt{+k} novo?
  \end{block}
  \begin{block}{Exercício 8 --- ordenar com uma pilha auxiliar}
    Ordene uma pilha (menor no topo) usando só uma segunda pilha e uma
    variável \texttt{int}. Quantos \texttt{empilha} isso faz no pior caso?
  \end{block}
  \begin{block}{Exercício 9 --- a lista de trás para a frente, sem recursão}
    Imprima uma lista simples do último ao primeiro com uma pilha de
    \texttt{No *}. Compare com a versão recursiva: quem faz o papel da
    pilha lá?
  \end{block}
  \fonte{Tenenbaum et al.\ (1995), cap.~2~\cite{tenenbaum1995estruturas};
    Cormen et al.\ (2009), §10.1, exercícios~\cite{cormen2009introduction}}
\end{frame}

\begin{frame}{Desafio}
  \begin{block}{Uma calculadora de verdade}
    \small Leia expressões \emph{infixas} com inteiros, \texttt{+ - * /} e
    parênteses, uma por linha, e imprima o valor:
    \vspace{-0.3em}
    \begin{itemize}
      \item primeiro converta para pós-fixa com uma pilha de
        \textbf{operadores}: um operador desempilha os de precedência maior
        ou igual antes de entrar; \texttt{(} entra sempre; \texttt{)}
        desempilha até o \texttt{(};
      \item depois avalie com uma pilha de \textbf{valores} (Exercício 4);
      \item parênteses desbalanceados e divisão por zero viram mensagem de
        erro, nunca \emph{segmentation fault}.
    \end{itemize}
  \end{block}
  \vspace{-0.3em}
  \centering\scriptsize
  \texttt{2 * (3 + 4) - 5} $\to$ \texttt{2 3 4 + * 5 -} $\to$ \texttt{9}
  \quad \texttt{(1 + 2} $\to$ \texttt{erro: parenteses}
  \fonte{Dijkstra (1961)~\cite{dijkstra1961algol};
    Tenenbaum et al.\ (1995), cap.~2~\cite{tenenbaum1995estruturas}}
\end{frame}

\begin{frame}{Fechamento}
  \begin{columns}[T]
    \begin{column}{0.5\textwidth}
      \textbf{O que mudou hoje}\\[0.3em]
      \small Restringir a estrutura a uma ponta deixou tudo mais simples:
      quatro operações, nenhum laço, e duas implementações trocáveis atrás
      de um \texttt{.h}.
    \end{column}
    \begin{column}{0.5\textwidth}
      \textbf{O que ficou}\\[0.3em]
      \small A pilha devolve na ordem \emph{inversa}. Para atender na ordem
      de chegada --- o primeiro a entrar, o primeiro a sair --- é preciso
      outra estrutura.
    \end{column}
  \end{columns}
  \begin{block}{Para levar}
    \begin{itemize}
      \small
      \item pilha: só o topo entra e sai --- o último a entrar sai primeiro;
      \item num vetor: topo em \texttt{v[n-1]}; confira vazia antes de
        tirar e cheia antes de pôr;
      \item vetor que cresce: dobre, não some um;
      \item numa lista: topo é o primeiro nó; a lista simples basta.
    \end{itemize}
  \end{block}
  \small\textcolor{UFSGray}{Próxima aula: filas --- põe-se por uma ponta e
  tira-se pela outra.}
\end{frame}

\section*{Referências}
\begin{frame}[allowframebreaks]{Referências}
  \nocite{*}
  \bibliographystyle{plain}
  \bibliography{referencias}
\end{frame}

\end{document}
